\chapter{Basic Commands}
\section{Defining the Basic Geometry of Calculation}
To define the basic geometry type of calculation, the command
\begin{verbatim} 
<geometry type="plane">
\end{verbatim} 
is used. Allowed parameters for {\bf type} are
{\bf 
\begin{itemize}
\item plane
\item axi
\item 3d
\end{itemize}}

If needed, a more detailed description of the geometry type must be
provided in the individual PDEs (e.g. in mechanics: {\bf "planeStrain"}).



\section{Defining the Input and Output File Format}
The format of the input and ouput files must be defined by the
commands {\bf input} and {\bf output}:

\begin{verbatim} 
<input  format="mesh" meshLibrary="cfsGrid"/>
<output format="gmv"/>
\end{verbatim} 

Allowed values of {\bf format} for input files are
{\bf \begin{itemize}
\item mesh
\item dat
\end{itemize}}
and for output files
{\bf \begin{itemize}
\item unv
\item gmv
\end{itemize}}
while the meshLibrary attribute can currently be one of
\begin{itemize}
\item \textbf{cfsGrid }
\item \textbf{adaptGrid}
\end{itemize}
\noindent
By default, the input format is \textbf{mesh}, the ouput format is \textbf{unv}
and the value for meshLibrary is \textbf{cfsGrid }.


\section{Defining the Material File}
To specify the name of the material file, the command 
\begin{verbatim}
<materialData file="matFileName"/>
\end{verbatim}
must be used. The default name for the material file is {\bf mat.dat}.

%
% =============================================================================
% TYPES OF ANALYSIS
% =============================================================================
%
\section{Analysis Techniques}\label{subsec:Analysis}
In the config-file, the command defining the analysis type reads as

\textbf{$<$analysis type=''TYPE''/$>$}

\noindent
where {\bf TYPE} can be chosen as one of
\begin{itemize} 
\item \xml{static} 
\item \xml{harmonic}
\item \xml{transient}
\item \xml{paramIdent}
\item \xml{multiSequence}
\item \xml{bubbleDynamic}
\item \xml{transient4Slice}
\item \xml{multiHarmonic}
\end{itemize}


% ----------------------------------------------------------------------------
\subsection{Static Analysis}
% ----------------------------------------------------------------------------
%
%
r static analysis, no additional information beside the command

{\bf \verb!<analysis type="static"/>!}

\noindent
is needed.

% ----------------------------------------------------------------------------
\subsection{Transient Analysis}
% ----------------------------------------------------------------------------
%
To be able to define the transient analysis in more detail, several
additional parameters have to be used:

\begin{verbatim}
<transient>
   <numSteps>    20    </numSteps>
   <firstDt>     5e-5  </firstDt>
   <stepSaveBeg> 1     </stepSaveBeg>
   <stepSaveEnd> 20    </stepSaveEnd>
   <stepSaveInc> 1     </stepSaveInc>
   <timeDataFile name="sin1kHz.dat"/>
   <writeRestartInc> 5 </writeRestartInc>
</transient>
\end{verbatim}

\noindent
Herein, the different tags have the following meaning:\\
\begin{tabular}{lll} 
$\circ$ & {\bf numSteps}    &  number of time steps\\
$\circ$ & {\bf firstDt }    &  time step size\\
$\circ$ & {\bf stepSaveBeg}  &  number of first step for saving results\\
$\circ$ & {\bf stepSaveEnd}  &  number of last step for saving results\\
$\circ$ & {\bf stepSaveIncr} & increment in steps for saving results\\
$\circ$ & {\bf timeDataFile} & file which holds the varying time values\\
$\circ$ & {\bf writeRestartInc} & step increment for writing a restart
file (currently only for mechanic PDE)
\end{tabular}
\bigskip

\noindent
Additionally, an input time-file must be provided:
\begin{verbatim} 
<timeDataFile name="sin1kHz.dat"/>
\end{verbatim} 
Therein, you have to define the time values in a column oriented
manner:
\begin{verbatim}
timeVal1          funcVal1
timeVal2          funcVal2
timeVal3          funcVal3
...
\end{verbatim}

% ----------------------------------------------------------------------------
\subsection{Harmonic Analysis}
% ----------------------------------------------------------------------------
%
Specification of parameters for a harmonic analysis is done by adding a
\xml{harmonic} section to the parameter file. When all currently available
parameters are set this section takes the following form:
%
%
\begin{verbatim}
  <harmonic>
    <numFreq>        ...   </numFreq>
    <startFreq>      ...   </startFreq>
    <stopFreq>       ...   </stopFreq>
    <sampling>       ...   </sampling>
    <adjustDamping>  ...   </adjustDamping>
    <matDataFreq>    ...   </matDataFreq>
  </harmonic>
\end{verbatim}
%
%
The parameters have the following meaning
%
%
\begin{center}
\begin{longtable}{llp{0.5\textwidth}}
\hline\noalign{\smallskip}
parameter & values & meaning \\
\noalign{\smallskip}\hline\noalign{\smallskip}
numFreq & $\nat$ & specifies the number $N$ of frequencies for which
                   a simulation is to be performed\\
\noalign{\smallskip}\hline\noalign{\smallskip}
startFreq & $> 0$ & sets $f_1$, i.e. the smallest frequency for which a
                    simulation is to be performed\\
\noalign{\smallskip}\hline\noalign{\smallskip}
stopFreq & $\geqslant$ startFreq & sets $f_N$, i.e. the largest frequency for
                                   which a simulation is to be performed\\
\noalign{\smallskip}\hline\noalign{\smallskip}
sampling & linear & in this case a linear sampling of the interval
                    $[f_1,f_N]$ is performed, i.e.
                    \begin{displaymath}
                    f_k = f_1 + (k-1) \frac{f_N - f_1}{N-1}
                    \enspace,\mbox{ for } k = 1,\ldots,N
                    \end{displaymath}\\
         & log    & in this case a logarithmic sampling of the interval
                    $[f_1,f_N]$ is performed, i.e.
                    \begin{displaymath}
                    f_k = f_1 \cdot
                    \left(\frac{f_N}{f_1}\right)^{\frac{k-1}{N-1}}
                    \enspace,\mbox{ for } k = 1,\ldots,N
                    \end{displaymath}
                    This leads to a finer sampling at the beginnig of the
                    interval and a coarser sampling at the end, since the
                    distance between samples increases as
                    \begin{displaymath}
                    f_{k+1} - f_k = \left[
                    \left(\frac{f_N}{f_1}\right)^{\frac{1}{N-1}}
                    - 1 \right] \cdot f_k
                    \end{displaymath}
                    and corresponds to a linear sampling of the transformed
                    interval $[\log_{10}(f_1),\log_{10}(f_N)]$.\\[1ex]
         & reverseLog & in this case a reverse logarithmic sampling of the
                        interval $[f_1,f_N]$ is performed, i.e.
                    \begin{displaymath}
                    f_k = f_N + f_1\cdot\left[ 1 -
                    \left(\frac{f_N}{f_1}\right)^{\frac{N-k}{N-1}}\right]
                    \end{displaymath}
                    for $k = 1,\ldots,N$.
                    This leads to a coarser sampling at the beginnig of the
                    interval and a finer sampling at the end.\\
\noalign{\smallskip}\hline\noalign{\smallskip}
adjustDamping & yes, no & set this to yes, if the damping is depending on the
                          frequency; in this case CFS++ will
                          adjust the damping for each frequency (for further
                          details ask Manfred \verb!;-)!\\
\noalign{\smallskip}\hline\noalign{\smallskip}
matDataFreq   & $>0$ & if adjustDamping is set to yes, this parameter
                       specifies the frequency of the damping parameters
                       in the material data file; if the parameter is omitted,
                       we assume that the damping parameters found are those
                       for startFreq\\
\noalign{\smallskip}\hline
\end{longtable}
\end{center}
%
%
% ----------------------------------------------------------------------------
\subsection{Multi--Sequence Analysis}
% ----------------------------------------------------------------------------
%
The idea of multi-sequence analysis is to allow the user to combine different
types of analyses in one program run. It might, e.g.~be sensible to combine a
static and a transient analysis in order to compute the pre--stress in a
mechanically clamped piezoelectric material, before starting a transient
computation, or to perform several static computations one after the other,
where different aspects, like e.g.~coils are sucessively added to the setup.

In order to perform a multi sequence run the following steps need to be
performed:
%
%
\begin{itemize}
\item The type of analysis in the \xml{analysis} element must be set to
      \xml{multiSequence}.
\item A multiSequence block like e.g.~the following must be inserted into
      the XML parameter file
\begin{verbatim}
<multiSequence>
  <step index="1">
    <pde type="electrostatic" analysis="static" refTag="elecStep1"/>
    <pde type="mechanic" analysis="static" refTag="mechStep1"/>
  </step>
  <step index="2">
    <pde type="electrostatic" analysis="transient" refTag="elecStep2"/>
    <pde type="mechanic" analysis="transient" refTag="mechStep2"/>
  </step>
</multiSequence>
\end{verbatim}
      The multiSequence driver will read this parameter block and perform
      an analysis for each \xml{step} element in the sequence of their
      occurence. The \xml{refTag} attribute specified in each \xml{pde}
      subelement is a symbolic name that should be unique. It is used to
      locate the appropriate boundary conditions or transient parameters
      etc.~for this step in the file. Since these parameters are indispensible
      this attribute and also the \xml{anlysis} attribute must be specified.
\item In order to specify for each step different boundary conditions or
      transient parameters the corresponding elements can be tagged with a
      \xml{tag} attribute. The parameters for a certain step of the
      multi-sequence analysis are found by comparing the \xml{tag} attribute
      with the respective \xml{refTag} attributes of the pdes of that step.

      In order to specify different parameters for different steps the element
      can be given several times and the \xml{tag} must be chosen accordingly.

      It is also possible to give the \xml{tag} attribute the special value
      \xml{anyTag}. This is also the default value, since the \xml{tag}
      attribute is optional. The \xml{anyTag} value indicates that the
      parameters are valid for all steps of the multi sequence analysis.

      A third possibility is to give the \xml{tag} attribute not a single
      value, but a list of different values, like
      e.g.~\xml{step1, step4, step5}. This can be used to 'turn on' this
      specific parameter setting for all PDEs in those steps with the
      corresponding tags.

      The following list gives an overview of all elements / parameter
      sections that currently can be \emph{tagged} in the way described above
      \begin{itemize}
      \item \xml{bcsAndLoads}
      \item \xml{transient}
      \item \xml{harmonic}
      \item \xml{bubbleDynamic}
      \item \xml{magnet}
      \item the \xml{bubbles} child element of the \xml{acoustic} PDE
      \item and all types of coils
      \end{itemize}
 
\end{itemize}
%
% ----------------------------------------------------------------------------
\subsection{ParamIdent - Analysis and Multiharmonics}
% ----------------------------------------------------------------------------
%
For the documentation of parameter identification in piezoelectric equations
see Chapter \ref{paramIdent}.The analysistype multiharmonics is not yet
implemented. Once it shall simulate the piezoelectric equations where the
material parameters are a function of frequency.
%
%
\section{List of PDEs}
Inside the config file command
\begin{verbatim}
<pdeList>
   List of PDEs ...
</pdeList>
\end{verbatim}
one has to specify the involved PDEs.  Some possible PDE names are
given in the following table: 

\begin{center}
  \begin{tabular}{l|l}
    {\it Type}      & {\it Description}  \\\hline
    {\bf acoustic}     & Acoustic PDE  \\
    {\bf electrostatic} & Electrostic PDE \\
    {\bf mechanic}     & Mechanics PDE\\    
  \end{tabular}
\end{center}

\noindent
As an example, a coupled electrostatic, mechanical simulation must be
defined as
\begin{verbatim}
<pdeList>
   <mechanic>
      Definitions of Mech-PDE ...
   </mechanic>

   <electrostatic>
      Definitions of Electrostatic-PDE ...
   </electrostatic>
</pdeList>
\end{verbatim}



\section{Geometrical Definitions}
In the top level section {\bf $<$domain$>$} one has to define all
geometrical entities, such as regions and their materials, boundary
nodes, save nodes, ...
\begin{verbatim}
<domain>
  <region name="cantilev" material="silizium"/>
  <region name="gap"      material="air"/>

  <nodes  name="xDof"/>
  <nodes  name="yDof"/>
  <nodes  name="yLoad"/>

  <elements name="mySaveElem"/>

  <directivityNodes>
    <distance radius="0.5"/>
    <distance radius="1.0"/>
    <center x="0.0" y="0.0" z="0.0"/>
    <planes xyName="myXYnodes" xyBegin="0" xyEnd="180"
            xzName="myXZnodes" xzBegin="0" xzEnd="360" 
            yzName="myYZnodes" yzBegin="-90" yzEnd="90"/>
    <saveIncr> 10 </saveIncr>
  </directivityNodes>

</domain>
\end{verbatim}

The keyword {\bf region} assigns a subdomain called {\bf name} to a material.
Further on, all in the simulation needed nodes (e.g. nodes which hold
a boundary condition) must be defined by the keyword {\bf nodes}.
Needed elements have to be defined by the keyword {\bf elements}.

Additionally, the optional element {\bf $<$directivityNodes$>$} is available if one wants to define (search) a collection of nodes near specific
radial distances from a center point (default 0, 0, 0) for later storing of results at
these nodes (e.g. for producing directivity plots in
acoustics). Nodes are selected on planes parallel to the principal planes. All 3D parameters are optional. XY plane is the default
plane for defining the nodes on a 2D simulation. Sweep angle for XY (xyEnd-xyBegin) is by default
360deg for the other planes it is 0deg. A name can be given to the nodes defined on each plane, otherwise default names will
also be used (dir\_XY, dir\_XZ, dir\_YZ). Finally, {\bf saveIncr}
represents the interval in degrees for defining the nodes.
\section{Numerical Integration Rules}
The numerical integration of the finite elements is done via Gaussian quadrature\footnote{See \emph{Higher-Order Finite Element Methods},
  Pavel Solin, Karel Segeth, Ivo Dolezel, 2004.}. Note that the integration rules are theoretically exact for the given order and hence a too low order of Ansatz-functions, a too coarse or deformed meshing cannot be compensated simply by a higher integration order. Table \ref{tab:default_order} gives the default order which should be sufficient for most simulations.

\begin{table}
\begin{center}
\begin{tabular}{r|r|r|r|r|r}
element & Ansatz-function & Gauss order & number of points \\
\hline
line & 1 & 1 & 1\\
     & 2 & 2=3 & 2 \\
\hline
triangle & 1 & 2 & 3 \\
         & 2 & 3 & 4 \\
\hline
rectangle & 1 & 2=3 & 4 \\
          & 2 & 2=3 & 4 \\
\hline
tetrahedron & 1 & 1 & 1 \\
         & 2 & 2 & 4 \\
\hline
hexahedron & 1 & 2=3 & 6 \\
         & 2 & 4=5 & 14 \\
\end{tabular}
\caption{This are the default settings for numerical integration of elements via Gaussian quadrature. Note that the integration rules for odd and even orders coincide for line, rectangle and hexahedron elements.}
\label{tab:default_order}
\end{center}
\end{table}
 
One can optionally define a different integration rule in the following way:
\begin{verbatim}
<integRules>
  <type>Gauss_standard</type>
  <order>5</order>
</integRules>
\end{verbatim}

The setting is applied on all element types, even if they are not used in the simulation. If the requested integration rule is not known for an element, CFS uses a default setting and prints a warning message stating all available integration rules for that element.

\begin{center}
\begin{tabular}{l|l}
integration type & description \\
\hline
Gauss\_standard & The so called economical Gaussian quadrature rules. \\
Classic & The legacy implementation, not of general interest. \\
Lobatto & An alternative integration rule for 1D. \\
Chebyshev & An alternative integration rule for 1D. \\
\end{tabular}
\end{center}

%%% Local Variables: 
%%% mode: latex
%%% TeX-master: "singlefields"
%%% End: 
